\section{État signalétique et des services}\label{annexe-services}

Le registre matricule du recrutement de Toulouse (classe 1896, n\textsuperscript{o}~911) est connu par une photographie conservée dans la famille. Les mentions en sont relevées ci-dessous dans l'ordre chronologique ; les abréviations du registre sont développées, les lectures douteuses signalées par un point d'interrogation. Les dates de passage dans les réserves, portées dans une grille dont les intitulés sont masqués par un cartouche collé, ne sont interprétées qu'avec réserve.

\begin{center}
\small
\renewcommand{\arraystretch}{1.18}
\begin{longtable}{@{}>{\raggedright\arraybackslash}p{27mm}>{\raggedright\arraybackslash}p{110mm}@{}}
\toprule
\multicolumn{2}{@{}l}{\textit{Signalement et état civil}}\\
\midrule
Signalement & Cheveux et sourcils bruns, yeux châtains, front découvert, nez pointu, bouche moyenne, menton rond, visage ovale ; taille 1,66~m. Degré d'instruction générale : 4.\\
Recensement & Classe 1896, n\textsuperscript{o}~911 au recrutement de Toulouse ; résidant à Toulouse, « étudiant en droit » ; parents domiciliés 47, rue du Taur, à Toulouse. N\textsuperscript{o}~74 du tirage au sort, canton de Toulouse-Nord.\\
Conseil de révision & « Exempté après ajournement de deux ans — faiblesse générale ».\\
Dates de passage & 1\textsuperscript{er} novembre 1900, 1910, 1916 et 1922 : passages successifs dans la réserve de l'armée active, l'armée territoriale et sa réserve, puis libération (interprétation probable).\\
\midrule
\multicolumn{2}{@{}l}{\textit{Services}}\\
\midrule
18 déc. 1914 & Classé service armé par le conseil de révision d'Eure-et-Loir (décret du 9 septembre 1914 sur la révision des exemptés).\\
20 mars 1915 & Rappelé à l'activité ; arrivé au corps, 10\textsuperscript{e} régiment d'artillerie, le 21 mars 1915. Campagne à l'intérieur du 21 mars au 3 mai 1915.\\
4 mai 1915 & Détaché à la Mission militaire française attachée à l'armée britannique, comme interprète (classe illisible : « 2\textsuperscript{e} cl. » ?) ; parti aux armées.\\
15 oct. 1915 & Passé au 19\textsuperscript{e} escadron du train des équipages militaires (décision ministérielle n\textsuperscript{o}~7239 3/1 du 3 octobre 1915).\\
4 mai 1915 -- 11 nov. 1918 & Campagne aux armées (campagne double) ; du 12 au 24 novembre 1918, campagne simple.\\
24 nov. 1918 & Dirigé sur l'École préparatoire des Andelys « comme professeur de langue anglaise ».\\
4 déc. 1918 & Nommé caporal.\\
17 (?) déc. 1918 & Passé au 10\textsuperscript{e} régiment d'artillerie à pied (circulaire ministérielle n\textsuperscript{o}~2332 1/11 du 8 janvier 1918).\\
1\textsuperscript{er} févr. 1919 & Passé au 19\textsuperscript{e} escadron du train.\\
7 févr. 1919 & Envoyé en congé illimité, à Nogent-le-Rotrou.\\
1919--1920 & Domiciles : Nogent-le-Rotrou, rue Saint-Hilaire (30 avril 1919) ; Montauban, 55, rue du Moustier, « prof. au lycée » (24 janvier 1920).\\
21 (?) avril 1923 & Affecté au 17\textsuperscript{e} escadron du train hippomobile (liste de sortie du 17\textsuperscript{e} corps d'armée n\textsuperscript{o}~1290 du 19 avril 1923).\\
10 nov. 1925 & Libéré du service militaire.\\
\bottomrule
\end{longtable}
\end{center}

La fiche ne dit rien de l'affectation d'Hippolyte Combes dans l'armée britannique. Sa lettre de février 1916 (ci-dessus, document~1) et le carnet la précisent : interprète à Saint-Omer à l'été de 1915, puis à la 18\textsuperscript{e} division, secrétaire pendant quinze jours au bureau divisionnaire de la Mission française, enfin, à partir d'août 1915, interprète de la 152\textsuperscript{nd} Company de l'\textit{Army Service Corps}, train de la division, jusqu'au 1\textsuperscript{er} novembre 1916 ; il est ensuite employé aux achats de paille et de fourrage avec le lieutenant Curtis, à Acheux puis à Hérissart. Le passage au train des équipages (octobre 1915) est un rattachement administratif, qui ne change rien à son service auprès des Britanniques. La tradition familiale d'un premier service au 1\textsuperscript{st} battalion des \textit{Irish Guards}, auprès du prince de Galles, n'est confirmée par aucune pièce.
